\ncertchapter{2}{Motion in a Straight Line}\chlabel{ch:straightline}

Kinematics is the description of motion without asking what causes it. That restriction can seem
artificial, since the interesting question is surely \emph{why} a thing moves, but it is the right
place to start: before a law of motion can be stated, we need a precise language in which to state it,
and that language is position, velocity and acceleration, together with the relations between them.

This chapter develops that language in one dimension, where every idea appears in its simplest form.
Nothing is lost by the restriction. When \chref{ch:vectors} introduces vectors, every definition here will carry
over unchanged, applied to each component in turn. The chapter has four parts: the distinction
between displacement and distance, which students routinely blur; the definitions of velocity and
acceleration as rates, together with the graphical meaning of slope and area; the three equations
of constant acceleration, derived rather than stated so that their one condition of validity,
constant acceleration, stays in view (forgetting it is the commonest error in kinematics); and the method for motion in which the acceleration varies.

\section{Position, Displacement and Path Length}

Throughout this chapter a moving body is treated as a \textbf{particle}: a point with no size. That is
a good model whenever the body is small compared with the distances it covers, or whenever every part of
it moves alike (a block sliding without turning); it fails for a spinning wheel, which waits for
\chref{ch:rotation}. To say where the particle is, fix an origin and a positive direction along the line,
together called a \textbf{frame of reference}; its \textbf{position} $x$ (in metres) is then a signed
number, positive on one side of the origin and negative on the other, and it changes with the time $t$
(in seconds). Over an interval, two different quantities describe how far it went.

\begin{defbox}{Displacement and path length}
\textbf{Displacement} $\Delta x=x_f-x_i$, the final position $x_f$ minus the initial position $x_i$, is
the signed change in position (the symbol $\Delta$ always means ``final minus initial''), the straight-line shortcut
from start to end, which can be positive, negative, or zero. \textbf{Path length} (distance) is the
total ground covered, never negative and never less than $|\Delta x|$. Figure~\ref{fig:displacement}
compares the two for a trip that doubles back.
\end{defbox}

\begin{figure}[htbp]
\centering
\tikzfig{kin_displacement}
\figcaption{fig:displacement}{Displacement is the arrow from start to end, $\Delta x=+\SI{3}{m}$; path length counts every metre travelled, out and back, \SI{9}{m}. On a straight line the two agree only for motion that never reverses.}
\end{figure}

A runner completing one lap of a track has substantial path length but \emph{zero} displacement, a
distinction that separates average speed (path length over time) from average velocity (displacement over
time). Over a round trip $\Delta x=0$ exactly, so average velocity is zero while average speed is
emphatically not. In general \textbf{average speed $\ge$ magnitude of average velocity}, with
equality only when the motion never reverses, because the path length is never less than $|\Delta x|$.
(Instantaneous speed, though, \emph{is} $|v|$ exactly: over an infinitesimal interval there is no room
to double back.)

\why{This split between displacement and path length is the first appearance of a structure that runs through all of
physics: a \emph{state} quantity that cares only about endpoints, versus a \emph{path} quantity that
depends on the route taken. The same distinction separates conservative from non-conservative forces
in \chref{ch:energy}, and internal energy from heat and work in thermodynamics. It shows up here first,
in the simplest possible setting.}

\section{Velocity and Acceleration}

Displacement says how far a particle moved, but not how quickly, and a sprinter and a snail can share the
same displacement. To compare them we need a rate: displacement per unit time. And because velocity
itself changes (a car pulling away, a ball slowing as it rises) we need a second rate, the change in
velocity per unit time. Velocity and acceleration are both rates, and each comes in two forms: an average
over an interval, and an instantaneous value obtained as that interval shrinks to nothing, the derivative
of \chref{ch:toolkit}.

\begin{defbox}{The kinematic quantities}
\textbf{Average velocity} $\bar v=\dfrac{\Delta x}{\Delta t}$; \textbf{instantaneous velocity}
$v=\dfrac{dx}{dt}$, the slope of the position--time graph. \textbf{Average acceleration}
$\bar a=\dfrac{\Delta v}{\Delta t}$; \textbf{instantaneous acceleration} $a=\dfrac{dv}{dt}=\dfrac{d^2x}{dt^2}$,
the slope of the velocity--time graph. Velocity is measured in \si{m/s} and acceleration in
\si{m/s^2} (dimensions $\mathsf{LT^{-1}}$ and $\mathsf{LT^{-2}}$). Both have a direction as well as a
size (they are vectors, \chref{ch:vectors}); in one dimension the direction is carried by the sign: $v>0$
means moving in the positive direction. The \textbf{speed} is
$|v|$, a scalar.
\end{defbox}

\begin{trapbox}{negative acceleration does not mean slowing down}
The sign of $a$ gives only its direction. A particle \textbf{speeds up} when $v$ and $a$ have the
\emph{same} sign and \textbf{slows down} (``retardation'') when they are opposite. Take $x=t^{3}-6t^{2}$
(metres, seconds), so $v=3t^{2}-12t$ and $a=6t-12$. At $t=\SI{1}{s}$, $v=\SI{-9}{m/s}$ and
$a=\SI{-6}{m/s^2}$: the acceleration is negative, yet the particle is speeding up. At $t=\SI{3}{s}$,
$v=\SI{-9}{m/s}$ and $a=\SI{+6}{m/s^2}$: it is slowing down. At $t=\SI{4}{s}$, $v=0$ but
$a=\SI{12}{m/s^2}$: momentarily at rest, still accelerating. Zero velocity never implies zero
acceleration.
\end{trapbox}

\subsection{Reading graphs}

Graphs carry kinematic information more compactly than equations, provided two operations are kept straight:
taking a slope, and finding an area. Figure~\ref{fig:vtgraph} shows both on a single velocity--time graph.

\begin{figure}[htbp]
\centering
\tikzfig{kin_vtgraph}
\figcaption{fig:vtgraph}{On a velocity--time graph, the \emph{slope} at a point is the acceleration and
the \emph{area} under the curve is the displacement. This dual relationship is the heart of graphical
kinematics.}
\end{figure}

\begin{keybox}{Slope and area}
\emph{Slope} turns position into velocity into acceleration (differentiate).
\emph{Area} turns acceleration into velocity into displacement (integrate). Position $\to$ velocity
$\to$ acceleration by slope; the reverse by area.
\end{keybox}

Why an area? Over a short interval $\Delta t$ the velocity barely changes, so the displacement is
$v\,\Delta t$, the area of a thin strip under the $v$--$t$ graph; adding the strips gives the whole area,
exactly the Toolkit's picture of an integral. The signed area is displacement; the \emph{unsigned} area is path length (area below the axis counts as
negative for displacement but positive for distance, exactly the walk-back case). Several things can
then be read off a graph with no formula at all:

\begin{itemize}
\item \textbf{Meeting:} two position--time curves crossing means the bodies are at the same place at the
same time.
\item \textbf{Same velocity at an instant:} the two position--time curves have equal slopes at that
time.
\item \textbf{At rest:} a horizontal position--time line. \textbf{Turning point:} a maximum or minimum,
where $v=0$ and is about to change sign.
\item \textbf{An $x$--$t$ graph can never be vertical}; that would be infinite velocity. If you have
drawn one, it is wrong.
\end{itemize}

\section{The Kinematic Equations}

Graphs show the relations between $x$, $v$ and $a$, but to get numbers we need equations. The simplest case,
and a surprisingly common one (free fall, a car braking steadily), is constant acceleration. Then the
definitions of velocity and acceleration can be integrated once and for
all, and the result is the three equations used in almost every problem of this chapter. They are derived
here rather than stated, so that their single condition of validity stays in view.

\begin{derivbox}{The three equations of constant acceleration}
Start from $a=\dfrac{dv}{dt}=$ const, with velocity $u$ and position $x_0$ at $t=0$. Inside an integral
the time variable is written $t'$, to keep it distinct from the upper limit $t$ (here the prime does not
mean a derivative).
\stepn{1} Integrate once: $\int_u^v dv=\int_0^t a\,dt'$ gives $v-u=at$, that is, $v=u+at$.
\stepn{2} Since $v=\dfrac{dx}{dt}$, integrate again: $\int_{x_0}^{x}dx=\int_0^t(u+at')\,dt'$ gives
$s=ut+\tfrac12 at^2$, where $s=x-x_0$ is the displacement from the starting position.
\stepn{3} Eliminate $t$ between the two: from the first, $t=(v-u)/a$ (for $a\neq0$), and substituting,
\[
s=u\frac{v-u}{a}+\frac{a}{2}\frac{(v-u)^2}{a^2}=\frac{v^2-u^2}{2a}
\;\Longrightarrow\;v^2=u^2+2as.
\]
These three, $v=u+at$, $s=ut+\tfrac12at^2$, $v^2=u^2+2as$, are the whole toolkit. One corollary is
worth having: substituting $a=(v-u)/t$ into step~2 gives $s=\tfrac12(u+v)\,t$, so under constant
acceleration the average velocity is the mean of the initial and final velocities. Setting $a=0$
recovers uniform motion, $s=ut$.
\end{derivbox}

There is a further result worth knowing: the distance covered in the $n$th \emph{second} of motion
(the interval from $t=n-1$ to $t=n$ seconds, $n$ a whole number),
\[
s_n=s(n)-s(n-1)=\big[un+\tfrac12an^2\big]-\big[u(n-1)+\tfrac12a(n-1)^2\big]=u+\tfrac{a}{2}(2n-1),
\]
where the right-hand side is understood to be multiplied by the one-second interval, which is why it
looks like a velocity yet is a distance. Derive it once and you never have to trust a memorised
version. And
note the third equation earns its place by containing no $t$: when a question mentions no time and asks
for none, $v^2=u^2+2as$ is the equation, and reaching for the others forces you to find $t$ only to throw
it away.

\inwords{$v^2=u^2+2as$ is the work--energy theorem in disguise. Multiply through by $\tfrac12 m$, for a body of mass
$m$, and anticipate Newton's second law, force $F=ma$ (\chref{ch:newton}):
$\tfrac12mv^2=\tfrac12mu^2+mas=\tfrac12mu^2+Fs$. \chref{ch:energy} derives this properly, for \emph{any}
force, but it was hiding in the kinematics all along.}

\begin{trapbox}{they hold only for constant acceleration}
The three equations assume $a$ is constant. They do \emph{not} apply when acceleration varies (as with\
air resistance, or a spring force). For those, return to $a=dv/dt$ and integrate, as in
\S\ref{sec:varacc}. Note also that $s$ is the \emph{displacement}, not the distance: if the body
reverses within the interval, split the motion at the turning point ($v=0$) and add the magnitudes.
\end{trapbox}

The commonest constant acceleration is \textbf{free fall}. Near the Earth's surface, with air
resistance negligible, every body, whatever its mass, accelerates downwards at the same
$g\approx\SI{9.8}{m/s^2}$ (often rounded to $10$), as Galileo found; why mass drops out is explained in
\chref{ch:newton} and \chref{ch:gravitation}. With upwards taken as positive, $a=-g$ throughout the
flight: going up, at the top and coming down.

\begin{exbox}{A ball thrown up}
A ball is thrown straight up at \SI{20}{m/s}. Take $g=\SI{10}{m/s^2}$ downward and neglect air
resistance. How high does it rise, and when does it return? Take upwards as positive, so $a=-g$. At the top $v=0$, and the time-free equation with $s=h$ gives
$v^2=u^2-2gh$: $0=20^2-2(10)h\Rightarrow h=\SI{20}{m}$.
Time up: $v=u-gt\Rightarrow0=20-10t\Rightarrow t=\SI{2}{s}$; by symmetry it returns \SI{2}{s} later, a
total of \SI{4}{s}. It lands with the speed it was thrown at: back at the start $s=0$, so
$v^2=u^2$ and $v=-\SI{20}{m/s}$, the minus sign meaning downwards. And after \SI{3}{s} the displacement is
$s=20(3)-\tfrac12(10)(3)^2=\SI{15}{m}$, but the distance travelled is $20+5=\SI{25}{m}$, because the ball turned
round at $t=\SI{2}{s}$ and fell $\tfrac12(10)(1)^2=\SI{5}{m}$ in the third second. \ans{\SI{20}{m}; \SI{4}{s}}
\end{exbox}

\section{When the Acceleration Is Not Constant}\label{sec:varacc}

Free fall and steady braking are idealisations. A rocket whose thrust changes, a body in air whose drag
grows with speed, or a mass on a spring all have an acceleration that varies, and for them the three
equations give wrong answers. Fortunately the definitions $v=dx/dt$ and $a=dv/dt$ assumed nothing about
the acceleration; only the three equations
did. When $a$ varies, we go back to the definitions and integrate them directly, using the Toolkit's
picture of an integral as a sum of thin slices (\chref{ch:toolkit}). Which integral to write depends on
what the acceleration is given as a function of.

\begin{derivbox}{Velocity and position from a varying acceleration}
\stepn{1} \textbf{Acceleration given as a function of time, $a(t)$.} Integrate $a=dv/dt$ and then
$v=dx/dt$ from the initial instant:
\[
v(t)=u+\int_0^t a(t')\,dt',\qquad x(t)=x_0+\int_0^t v(t')\,dt'.
\]
Graphically, the change in velocity is the area under the $a$--$t$ graph.
\stepn{2} \textbf{Acceleration given as a function of position, $a(x)$.} Time is not available to
integrate over, so remove it with the chain rule:
$a=\dfrac{dv}{dt}=\dfrac{dv}{dx}\,\dfrac{dx}{dt}=v\dfrac{dv}{dx}$. Then $v\,dv=a\,dx$, and integrating,
\[
\tfrac12v^2-\tfrac12u^2=\int_{x_0}^{x}a(x')\,dx'.
\]
\stepn{3} \textbf{Check.} For constant $a$ the integral in Step~2 is $a(x-x_0)=as$, and the result
becomes $v^2=u^2+2as$, as it must.
\end{derivbox}

\begin{exbox}{A thrust that dies away}
A particle starts from rest at the origin with acceleration $a=6-2t$ (in \si{m/s^2}, with $t$ in
seconds). Find its maximum velocity and its position at that moment.

\stepn{1} \textbf{Velocity.} $v=\int_0^t(6-2t')\,dt'=6t-t^2$.
\stepn{2} \textbf{Maximum.} $v$ is greatest where $dv/dt=a=0$, that is at $t=\SI{3}{s}$, giving
$v=18-9=\SI{9}{m/s}$. After that the acceleration is negative and the particle slows.
\stepn{3} \textbf{Position.} $x=\int_0^t(6t'-t'^2)\,dt'=3t^2-\tfrac13t^3$, so at $t=\SI{3}{s}$,
$x=27-9=\SI{18}{m}$.

Had we used the initial acceleration $a=\SI{6}{m/s^2}$ in the constant-acceleration equations, we would
have found \SI{18}{m/s} and \SI{27}{m}: wrong by a factor of two in the velocity. \ans{\SI{9}{m/s} at
$x=\SI{18}{m}$}
\end{exbox}

\section{Summary}\label{sec:summary2}
Displacement is a signed change in position; distance is the total ground covered and is never smaller
than the magnitude of the displacement. Velocity is the rate of change of position and acceleration the
rate of change of velocity, and on a graph each rate is a slope while each accumulation is an area.
When the acceleration is constant, integrating twice gives the three kinematic equations, and choosing
among them is a matter of noticing which of the five quantities $u$, $v$, $a$, $s$, $t$ the question
neither gives nor asks for. When the acceleration varies, the equations fail and the definitions must
be integrated directly, using $a=v\,dv/dx$ if the acceleration is given as a function of position.

\begin{itemize}
\item $v=dx/dt$ and $a=dv/dt$: the slope of $x$--$t$ is velocity, the slope of $v$--$t$ is acceleration.
\item The area under $v$--$t$ is displacement; the area under $a$--$t$ is change in velocity.
\item $v=u+at$, $\;s=ut+\tfrac12at^2$, $\;v^2=u^2+2as$, valid \emph{only} for constant $a$.
\item Distance in the $n$th second: $s_n=u+\tfrac{a}{2}(2n-1)$, which is $s(n)-s(n-1)$ and should be
      derived, not recalled.
\item At the top of a vertical throw $v=0$ but $a=-g$: zero velocity never implies zero acceleration.
\item $v^2=u^2+2as$ is the work--energy theorem in disguise; \chref{ch:energy} makes that precise.
\end{itemize}

\section{Problem Types}\label{sec:types2}
Straight-line kinematics is one of the largest sources of examination questions in mechanics, and the
six types below cover the bulk of them. In every one, the decisive first step is the same: choose a positive direction and
hold to it, because nearly every wrong answer in this chapter is a sign error rather than a conceptual one.

\archq{Constant acceleration: choosing the equation}
{A body moves with constant acceleration; find a distance, time, or velocity.}
{Pick the equation that omits the quantity you neither know nor want. List $u,v,a,s,t$: three are
known, one is wanted, and the remaining one is absent, which names the equation.}
{Sign conventions. Choose one positive direction and stick to it; upward-thrown motion has $a=-g$
throughout, including at the top where $v=0$ but $a\ne0$.}
{A car brakes uniformly from \SI{30}{m/s} to rest in \SI{5}{s}. Its stopping distance is:}
{\SI{37.5}{m}}{\SI{75}{m}}{\SI{150}{m}}{\SI{225}{m}\archans{B}{No acceleration is given or asked, so
use $s=\tfrac12(u+v)t=\tfrac12(30+0)(5)=\SI{75}{m}$. Option (C) is $ut$, as though the car never slowed;
option (D) adds $\tfrac12at^2$ with the wrong sign of $a$.}}

\archq{Distance travelled in the $n$th second}
{``How far does the body travel in the 5th second?''; note: \emph{in} the $n$th second, not \emph{in}
$n$ seconds.}
{Use $s_n=u+\dfrac{a}{2}(2n-1)$. It is $s(n)-s(n-1)$, so despite its look it is a
\textbf{distance}, measured in metres, not a velocity.}
{Substituting into $s=ut+\tfrac12at^2$ with $t=n$, which gives the distance in the first $n$ seconds,
a much larger number.}
{A body starts from rest with $a=\SI{2}{m/s^2}$. The distance covered in the 5th second is:}
{\SI{9}{m}}{\SI{4}{m}}{\SI{25}{m}}{\SI{50}{m}\archans{A}{$s_5=0+\tfrac22(2\times5-1)
=1\times9=\SI{9}{m}$. Option (C) is $s(5)=\tfrac12(2)(25)$, the distance in five seconds, the classic
misread.}}

\archq{Stopping distance and its scaling}
{``A vehicle braking with the same retardation from double the speed, how does the stopping distance
change?''}
{Use the time-free equation: $0=u^2-2as$ gives $s=\dfrac{u^2}{2a}$, so $\boxed{s\propto u^2}$ for a fixed
retardation. Double the speed, \emph{quadruple} the distance.}
{Assuming the distance doubles when the speed doubles. The dependence is quadratic, which is exactly the
point of the question.}
{A car travelling at \SI{20}{m/s} stops in \SI{20}{m}. Under the same braking, its stopping distance
from \SI{40}{m/s} is:}
{\SI{40}{m}}{\SI{60}{m}}{\SI{160}{m}}{\SI{80}{m}\archans{D}{$s\propto u^2$, and the speed is doubled,
so $s$ is multiplied by $2^2=4$: $4\times20=\SI{80}{m}$. Option (A) is the linear guess.}}

\archq{Reading a motion graph}
{``From the given $x$--$t$ or $v$--$t$ graph, find the velocity, acceleration or displacement.''}
{On an $x$--$t$ graph the \emph{slope} is velocity. On a $v$--$t$ graph the \emph{slope} is acceleration
and the \emph{area} is displacement. Area below the axis counts as \emph{negative} displacement, so
for total \emph{distance}, add the magnitudes instead.}
{Confusing distance with displacement when the graph crosses the axis, and reading a curved $x$--$t$
graph as though constant velocity.}
{A body's $v$--$t$ graph is a triangle: $v$ rises linearly from $0$ to \SI{20}{m/s} in \SI{5}{s}, then
falls linearly back to $0$ in another \SI{5}{s}. The total displacement is:}
{\SI{50}{m}}{\SI{100}{m}}{\SI{200}{m}}{\SI{400}{m}\archans{B}{Displacement is the area of the triangle:
$\tfrac12\times\text{base}\times\text{height}=\tfrac12(10)(20)=\SI{100}{m}$. Option (C) forgets the
$\tfrac12$, treating the area as a rectangle.}}

\archq{Free fall and vertical throw}
{``A body is dropped/thrown up, find the time to fall, the maximum height, or the speed on
return.''}
{Take one sign convention and hold it. For a body thrown up at $u$: rise time $=u/g$, maximum height
$=u^2/2g$, total flight $=2u/g$, and, with air resistance neglected, it \textbf{returns with the same
speed} $u$. For a dropped body
$h=\tfrac12gt^2$.}
{Forgetting that the upward and downward journeys are symmetric, or changing the sign of $g$ halfway
through a problem.}
{A stone is dropped from rest. The distance it falls in \SI{3}{s} is ($g=\SI{10}{m/s^2}$):}
{\SI{15}{m}}{\SI{30}{m}}{\SI{45}{m}}{\SI{90}{m}\archans{C}{$h=\tfrac12gt^2=\tfrac12(10)(9)=\SI{45}{m}$.
Option (B) is the trap of computing $gt$, a velocity, not a distance.}}

\archq{Non-uniform acceleration, integrate instead}
{``Given $v(t)$ or $a(t)$ as a \emph{function}, find the displacement or velocity.'' The tell is that
$a$ is \emph{not} constant.}
{The constant-acceleration equations do \textbf{not} apply. Go back to the definitions and integrate:
$v=\int a\,dt$ and $s=\int v\,dt$, fixing the constants from the initial conditions. If $a$ is given as a
function of $x$, use the chain rule,
$a=\dfrac{dv}{dt}=\dfrac{dv}{dx}\dfrac{dx}{dt}=v\dfrac{dv}{dx}$.}
{Using $s=ut+\tfrac12at^2$ when $a$ depends on $t$. Those formulas were \emph{derived} assuming constant
acceleration and are invalid here.}
{A particle moves with $v=3t^2$ (with $v$ in \si{m/s} and $t$ in \si{s}) starting from $x=0$. Its displacement in the first \SI{2}{s} is:}
{\SI{4}{m}}{\SI{6}{m}}{\SI{8}{m}}{\SI{12}{m}\archans{C}{$s=\int_0^2 3t^2\,dt=\big[t^3\big]_0^2
=\SI{8}{m}$. Note $a=6t$ is \emph{not} constant, so no constant-acceleration formula applies.}}
